% GENERATED FILE. Edit canonical lessons, then run npm run generate:books.
% canonical-source-hash: fnv1a64:53ddb5348c5a48eb
% canonical-lessons: TA-C182-kulircatanap-petti, TA-C182-kulay, TA-C182-tamlar, TA-C182-catti, TA-C182-muti

\chapter{Refrigerator, Tap, Tumbler, Pan, Lid}
\label{ch:ta-refrigerator-tap-tumbler-pan-lid}
\hlchaptermodality{182}

\begin{chapteropening}
\textbf{By the end of this chapter:} \emph{I can say a refrigerator, a tap, a tumbler, a pan and a lid.}

Complete the last lesson of chapter 182: I can say a refrigerator, a tap, a tumbler, a pan and a lid.
\end{chapteropening}

\section[\texorpdfstring{\ta{குளிர்சாதனப்} \ta{பெட்டி} (kuḷircātaṉap peṭṭi)}{குளிர்சாதனப் பெட்டி (kuḷircātaṉap peṭṭi)}]{\ta{குளிர்சாதனப்} \ta{பெட்டி} (kuḷircātaṉap peṭṭi) — a refrigerator}

\label{lesson:TA-C182-kulircatanap-petti}



\begin{quote}
Before the new one: say the Tamil for everything, then the Tamil for thick.
\end{quote}

\subsection*{You'll want to know: \ta{குளிர்சாதனப்} \ta{பெட்டி}}
\textbf{\ta{குளிர்சாதனப்} \ta{பெட்டி}} — \emph{kuḷircātaṉap peṭṭi} — \textquotedblleft{}a refrigerator\textquotedblright{}.

\begin{grammarlens}[title={What you've built}]
One more word for this chapter.
\end{grammarlens}

\subsection*{Your turn}
\begin{itemize}
\raggedright
  \item \emph{Say it:} \emph{kuḷircātaṉap peṭṭi}
  \item \emph{Say it:} \emph{kuḷircātaṉap peṭṭi}, once more
  \item \emph{Say it:} say \emph{kuḷircātaṉap peṭṭi}
  \item \emph{Recall:} say the Tamil for everything, then the Tamil for thick, then say \emph{kuḷircātaṉap peṭṭi} again
\end{itemize}

\begin{tcolorbox}[breakable,colback=teal!4,colframe=teal!35!black,title={Before you move on}]
What does \ta{குளிர்சாதனப்} \ta{பெட்டி} mean? (\textquotedblleft{}A refrigerator\textquotedblright{}.) Say it once more.
\end{tcolorbox}
\section[\texorpdfstring{\ta{குழாய்} (kuḻāy)}{குழாய் (kuḻāy)}]{\ta{குழாய்} (kuḻāy) — a tap, a pipe}

\label{lesson:TA-C182-kulay}



\begin{quote}
Before the new one: say the Tamil for thick, then the Tamil for a refrigerator.
\end{quote}

\subsection*{You'll want to know: \ta{குழாய்}}
\textbf{\ta{குழாய்}} — \emph{kuḻāy} — \textquotedblleft{}a tap, a pipe\textquotedblright{}.

\begin{grammarlens}[title={What you've built}]
One more word for this chapter.
\end{grammarlens}

\subsection*{Your turn}
\begin{itemize}
\raggedright
  \item \emph{Say it:} \emph{kuḻāy}
  \item \emph{Say it:} \emph{kuḻāy}, once more
  \item \emph{Say it:} say \emph{kuḻāy}
  \item \emph{Recall:} say the Tamil for thick, then the Tamil for a refrigerator, then say \emph{kuḻāy} again
\end{itemize}

\begin{tcolorbox}[breakable,colback=teal!4,colframe=teal!35!black,title={Before you move on}]
What does \ta{குழாய்} mean? (\textquotedblleft{}A tap, a pipe\textquotedblright{}.) Say it once more.
\end{tcolorbox}
\section[\texorpdfstring{\ta{டம்ளர்} (ṭamḷar)}{டம்ளர் (ṭamḷar)}]{\ta{டம்ளர்} (ṭamḷar) — a tumbler, a drinking glass}

\label{lesson:TA-C182-tamlar}



\begin{quote}
Before the new one: say the Tamil for a refrigerator, then the Tamil for a tap.
\end{quote}

\subsection*{You'll want to know: \ta{டம்ளர்}}
\textbf{\ta{டம்ளர்}} — \emph{ṭamḷar} — \textquotedblleft{}a tumbler, a drinking glass\textquotedblright{}.

\begin{grammarlens}[title={What you've built}]
One more word for this chapter.
\end{grammarlens}

\subsection*{Your turn}
\begin{itemize}
\raggedright
  \item \emph{Say it:} \emph{ṭamḷar}
  \item \emph{Say it:} \emph{ṭamḷar}, once more
  \item \emph{Say it:} say \emph{ṭamḷar}
  \item \emph{Recall:} say the Tamil for a refrigerator, then the Tamil for a tap, then say \emph{ṭamḷar} again
\end{itemize}

\begin{tcolorbox}[breakable,colback=teal!4,colframe=teal!35!black,title={Before you move on}]
What does \ta{டம்ளர்} mean? (\textquotedblleft{}A tumbler, a drinking glass\textquotedblright{}.) Say it once more.
\end{tcolorbox}
\section[\texorpdfstring{\ta{சட்டி} (caṭṭi)}{சட்டி (caṭṭi)}]{\ta{சட்டி} (caṭṭi) — a pan, a cooking pot}

\label{lesson:TA-C182-catti}



\begin{quote}
Before the new one: say the Tamil for a tap, then the Tamil for a tumbler.
\end{quote}

\subsection*{You'll want to know: \ta{சட்டி}}
\textbf{\ta{சட்டி}} — \emph{caṭṭi} — \textquotedblleft{}a pan, a cooking pot\textquotedblright{}.

\begin{grammarlens}[title={What you've built}]
One more word for this chapter.
\end{grammarlens}

\subsection*{Your turn}
\begin{itemize}
\raggedright
  \item \emph{Say it:} \emph{caṭṭi}
  \item \emph{Say it:} \emph{caṭṭi}, once more
  \item \emph{Say it:} say \emph{caṭṭi}
  \item \emph{Recall:} say the Tamil for a tap, then the Tamil for a tumbler, then say \emph{caṭṭi} again
\end{itemize}

\begin{tcolorbox}[breakable,colback=teal!4,colframe=teal!35!black,title={Before you move on}]
What does \ta{சட்டி} mean? (\textquotedblleft{}A pan, a cooking pot\textquotedblright{}.) Say it once more.
\end{tcolorbox}
\section[\texorpdfstring{\ta{மூடி} (mūṭi)}{மூடி (mūṭi)}]{\ta{மூடி} (mūṭi) — a lid}

\label{lesson:TA-C182-muti}



\begin{quote}
Before the new one: say the Tamil for a tumbler, then the Tamil for a pan.
\end{quote}

\subsection*{You'll want to know: \ta{மூடி}}
\textbf{\ta{மூடி}} — \emph{mūṭi} — \textquotedblleft{}a lid\textquotedblright{}.

\begin{grammarlens}[title={What you've built}]
One more word for this chapter.
\end{grammarlens}

\subsection*{Your turn}
\begin{itemize}
\raggedright
  \item \emph{Say it:} \emph{mūṭi}
  \item \emph{Say it:} \emph{mūṭi}, once more
  \item \emph{Say it:} say \emph{mūṭi}
  \item \emph{Recall:} say the Tamil for a tumbler, then the Tamil for a pan, then say \emph{mūṭi} again
\end{itemize}

\begin{tcolorbox}[breakable,colback=teal!4,colframe=teal!35!black,title={Before you move on}]
What does \ta{மூடி} mean? (\textquotedblleft{}A lid\textquotedblright{}.) Say it once more.
\end{tcolorbox}
